\section{Objetivos}

\subsection{Objetivo general}

Implementar en Go, sin librerías de terceros, una versión secuencial y una concurrente del mismo
algoritmo K-means sobre más de un millón de viajes reales; medir con rigor estadístico cuánto
acelera la concurrencia y a qué costo; y verificar formalmente que su sincronización está libre de
deadlocks y respeta la exclusión mutua, contrastando además el código con un análisis asistido por
IA.

\subsection{Objetivos específicos}

\begin{itemize}
  \item Conservar los puntos de los Entregables 1 y 2 y corregir de forma expresa lo que cambió
        desde cada entrega.
  \item Plantear el algoritmo concurrente y modelar su sincronización en Promela, de modo que Spin
        recorra todos los entrelazados y confirme que ningún punto se pierde ni se duplica y que los
        centroides no se escriben mientras se leen.
  \item Implementar las dos versiones en Go con un contrato numérico común (misma inicialización,
        misma distancia, mismo criterio de parada) y probar su equivalencia con \texttt{go test -race}.
  \item Explicar el patrón \emph{worker pool} elegido y los mecanismos de sincronización que usa:
        canal de bloques, acumuladores privados y barrera con \texttt{sync.WaitGroup}.
  \item Calcular el speedup $T_{\text{secuencial}} / T_{\text{concurrente}}$ con media recortada
        sobre múltiples ejecuciones, con tabla, dispersión e intervalos de confianza.
  \item Analizar el speedup, la escalabilidad fuerte y débil, y los trade-offs entre la versión
        secuencial y la concurrente.
  \item Medir el uso de memoria y CPU de cada configuración y ubicar el punto de equilibrio a partir
        del cual la concurrencia paga.
  \item Verificar formalmente con Spin la ausencia de deadlocks, la exclusión mutua entre el
        coordinador y los workers, y la terminación bajo \emph{weak fairness}, con un mutante por
        propiedad que demuestre que cada verificación detecta el error que dice detectar.
  \item Analizar el código del repositorio con un modelo de IA y un prompt estructurado, y
        documentar en Markdown los GAPs en calidad de código, seguridad y patrones de concurrencia.
  \item Presentar las conclusiones y recomendaciones de cada integrante.
  \item Evidenciar el trabajo en GitHub con el flujo Git Flow y la participación de los integrantes.
\end{itemize}
